Deux haut-parleurs branchés sur le même amplificateur émettent en phase le même son pur. Un auditeur se trouve à \qty{\v{r1}}{m} du premier haut-parleur et à \qty{\v{r2}}{m} du second. La vitesse du son est de \qty{\v{v}}{m/s}.

\begin{enumerate}[label=\alph*)]
    \item Déterminez la différence de marche entre les deux ondes au point où se trouve l'auditeur.
    \item Quelle est la plus basse fréquence pour laquelle l'auditeur entend une interférence destructive ?
    \item Quelle est la plus basse fréquence (non nulle) pour laquelle il entend une interférence constructive ?
    \item On inverse les fils d'un des haut-parleurs, de sorte que les sources sont maintenant déphasées de $\pi$. Quelle est alors la plus basse fréquence donnant une interférence constructive ?
\end{enumerate}
